% !TEX root = ../main.tex
\section{Understanding TSGD: Two Error Sources}
\label{sec:two-error-sources}

\subsection{A Running Quadratic Model}
\label{subsec:quadratic-model}

We use one scalar model throughout this section:
\begin{equation}
\label{eq:quadratic-noise-model}
f(x)=\frac{\lambda}{2}x^2,\qquad
h(x;a,b)=(\lambda+a)x+b,
\qquad \lambda>0.
\end{equation}
The noises $a$ and $b$ are independent, with
\[
\mathbb{E}[a]=\mathbb{E}[b]=0,\qquad
\mathbb{E}[a^2]=\delta^2,\qquad
\mathbb{E}[b^2]=\sigma^2.
\]
The objective has gradient $\nabla f(x)=\lambda x$, and
$h(x;a,b)$ is an unbiased stochastic estimate of it:
$\mathbb{E}[h(x;a,b)]=\nabla f(x)$.
At each transport step, a fresh pair $(a_{k+1},b_{k+1})$ is drawn
independently of the past and of the square-integrable initial estimate
$g_0$; the same pair is used at both endpoints. We take $x_0$ to be
deterministic and $\eta>0$ constant.

Although the pointwise gradient noise is $ax+b$, the additive term
disappears from a coupled difference:
\begin{equation}
\label{eq:quadratic-coupled-difference}
D_{k+1}=(\lambda+a_{k+1})(x_{k+1}-x_k).
\end{equation}
Indeed, since
$D_{k+1}-[\nabla f(x_{k+1})-\nabla f(x_k)]
=a_{k+1}(x_{k+1}-x_k)$,
\[
\mathbb{E}\bigl[|D_{k+1}-[\nabla f(x_{k+1})-
\nabla f(x_k)]|^2\mid\mathcal F_k\bigr]
=\delta^2|x_{k+1}-x_k|^2.
\]
Thus Assumption~\ref{ass:centered-similarity} holds with equality
for $\delta=\sqrt{\mathbb{E}[a^2]}$, independently of $\sigma^2$.
Using $x_{k+1}-x_k=-\eta g_k$ in the TSGD recursion gives,
with $e_k=g_k-\lambda x_k$,
\begin{equation}
\label{eq:quadratic-core-recursions}
g_{k+1}=(1-\eta\lambda-\eta a_{k+1})g_k,
\qquad
e_{k+1}=e_k-\eta a_{k+1}g_k.
\end{equation}
The first equation describes the estimator; the second shows how
transport changes its error. We retain $E_k$ from Section~2 and
$M_k,G_k$ from Section~3; in this scalar model,
$E_k=\mathbb{E}[|e_k|^2]$, $M_k=\mathbb{E}[|g_k|^2]$, and
$G_k=\mathbb{E}[|\nabla f(x_k)|^2]$. Thus $G_k$ measures the true
gradient at $x_k$, while $E_0$ and $M_0$ quantify the initial
estimation error and estimator magnitude, respectively.

\subsection{Exact Initialization and Exact Transport}
\label{subsec:exact-init-exact-transport}

First, the conclusion is not specific to the quadratic. Suppose
$g_0=\nabla f(x_0)$ and every transported difference is exact:
\[
D_{k+1}=\nabla f(x_{k+1})-\nabla f(x_k).
\]
If $g_k=\nabla f(x_k)$, the TSGD estimator update gives
$g_{k+1}=g_k+D_{k+1}=\nabla f(x_{k+1})$. Since this holds at
$k=0$, induction yields $g_k=\nabla f(x_k)$ for every $k$, and hence
\[
x_{k+1}=x_k-\eta\nabla f(x_k).
\]
Thus TSGD is exactly gradient descent, not merely equivalent to it
in expectation.

Under Assumption~\ref{ass:smooth-objective}, choose the standard GD
stepsize $\eta=1/L$. Theorem~\ref{thm:tsgd-stationarity} with
$E_0=\delta=0$ then recovers the nonconvex GD guarantee:
\begin{equation}
\label{eq:exact-transport-gd-bound}
\frac{1}{K}\sum_{k=0}^{K-1}\|\nabla f(x_k)\|^2
\leq\frac{2L\Delta_0}{K}.
\end{equation}
There is also a stronger pointwise conclusion at this fixed
stepsize. The GD descent inequality implies
\[
\frac{1}{2L}\sum_{k=0}^{K-1}\|\nabla f(x_k)\|^2
\leq f(x_0)-f(x_K)\leq\Delta_0.
\]
This bound holds for every $K$. The nonnegative partial sums are
therefore uniformly bounded, and letting $K\to\infty$ gives
$\sum_{k=0}^{\infty}\|\nabla f(x_k)\|^2\leq2L\Delta_0<\infty$.
Hence $\|\nabla f(x_k)\|^2\to0$; no convexity or convergence of the
iterates is needed.

In the running quadratic, exact initialization means
$g_0=\lambda x_0=\nabla f(x_0)$. Exact transport corresponds to
$\delta=0$, so $a=0$ almost surely, and the shared additive noise
$b_{k+1}$ cancels:
\[
D_{k+1}=\lambda(x_{k+1}-x_k)
=\nabla f(x_{k+1})-\nabla f(x_k).
\]
Consequently, $g_k=\lambda x_k=\nabla f(x_k)$ for every $k$.
Writing $r:=1-\eta\lambda$, the exact trajectory is
\begin{equation}
\label{eq:quadratic-exact-gd}
x_k=r^k x_0,\qquad
g_k=\lambda r^k x_0,\qquad E_k=0.
\end{equation}
For $0<\eta<2/\lambda$, we have $|r|<1$ and
$G_k=\lambda^2r^{2k}x_0^2\to0$.
The lesson is that correct initialization and exact transport recover
deterministic optimization despite large additive oracle noise.

\subsection{Inexact Initialization and Exact Transport}
\label{subsec:inexact-init-exact-transport}

Let $g_0=\nabla f(x_0)+e_0$ and suppose all transported differences
are exact. Then $e_k=e_0$ for every $k$, and hence
\[
g_k=\nabla f(x_k)+e_0
=\nabla\widetilde f_{e_0}(x_k),\qquad
\widetilde f_{e_0}(x):=f(x)+\langle e_0,x\rangle.
\]
Thus $E_k=E_0$ and TSGD is GD on the fixed, linearly tilted
objective $\widetilde f_{e_0}$. Exact transport cannot correct the
initial error.

Theorem~\ref{thm:tsgd-stationarity} with $\delta=0$ gives, at the
GD stepsize $\eta=1/L$,
\begin{equation}
\label{eq:exact-transport-inexact-init-bound}
\frac1K\sum_{k=0}^{K-1}G_k
\leq\frac{2L\Delta_0}{K}+E_0.
\end{equation}
If $e_0$ is fixed and $\widetilde f_{e_0}$ is bounded below,
GD on that objective also gives
$\nabla f(x_k)\to-e_0$ and
$\|\nabla f(x_k)\|^2\to\|e_0\|^2$.
The lower bound on $\widetilde f_{e_0}$ is needed only for this
pathwise asymptotic claim, not for \eqref{eq:exact-transport-inexact-init-bound}.

The obstruction is informational: for
$f_c(x)=f(x)+\langle c,x\rangle$,
$\nabla f_c(x)-\nabla f_c(y)=\nabla f(x)-\nabla f(y)$.
Gradient differences cannot identify an unknown constant offset $c$.

In the quadratic model, $\delta=0$ makes transport exact regardless
of the additive noise $b$. With $g_0=\lambda x_0+e_0$ and
$r=1-\eta\lambda$, the exact trajectory is
\begin{equation}
\label{eq:quadratic-tilted-trajectory}
g_k=r^k(\lambda x_0+e_0),\qquad
x_k=r^k x_0-\frac{1-r^k}{\lambda}e_0.
\end{equation}
If the initialization error is centered, $\mathbb E[e_0]=0$, then
\begin{equation}
\label{eq:quadratic-centered-init-stationarity}
G_k=r^{2k}\lambda^2x_0^2+(1-r^k)^2E_0.
\end{equation}
If $0<\eta<2/\lambda$, then $x_k\to-e_0/\lambda$ in mean
square and, even for biased $e_0$,
\begin{equation}
\label{eq:quadratic-init-residual}
\lim_{k\to\infty}G_k=E_0.
\end{equation}
For $\varepsilon>0$, choose $N$ such that
$|G_k-E_0|\leq\varepsilon$ for $k\geq N$. Then, for $K>N$,
\[
\left|\frac1K\sum_{k=0}^{K-1}G_k-E_0\right|
\leq\frac1K\sum_{k=0}^{N-1}|G_k-E_0|+\varepsilon.
\]
Letting $K\to\infty$ and then $\varepsilon\downarrow0$ gives,
for the admissible choice $\eta=1/L=1/\lambda$,
\[
\lim_{K\to\infty}\frac1K\sum_{k=0}^{K-1}G_k
=E_0
=\lim_{K\to\infty}
\left(\frac{2\Delta_0}{\eta K}+E_0\right).
\]
Thus the coefficient one on $E_0$ in
\eqref{eq:tsgd-stationarity-bound} is sharp.

\subsection{Exact Initialization and Stochastic Transport}
\label{subsec:exact-init-stochastic-transport}

Now take $g_0=\lambda x_0$ and $\delta>0$. Define
\[
r:=1-\eta\lambda,\qquad
\rho:=\mathbb{E}[(r-\eta a)^2]
=r^2+\eta^2\delta^2
=1-2\eta\lambda+\eta^2(\lambda^2+\delta^2).
\]
Independence of the fresh $a_{k+1}$ from the past and
\eqref{eq:quadratic-core-recursions} give the exact moment dynamics
\begin{equation}
\label{eq:quadratic-moment-recursions}
M_{k+1}=\rho M_k,\qquad
E_{k+1}=E_k+\eta^2\delta^2M_k.
\end{equation}
Since $E_0=0$ and $M_0=\lambda^2x_0^2$, these solve to
\begin{equation}
\label{eq:quadratic-exact-init-moments}
M_k=\rho^k\lambda^2x_0^2,\qquad
E_k=\eta^2\delta^2\lambda^2x_0^2
       \sum_{t=0}^{k-1}\rho^t.
\end{equation}
The estimator error is a martingale, as in Section~2.3, and its
mean-square size increases by the variance of each new transport
error.

The mean-square stability condition is
\begin{equation}
\label{eq:quadratic-stability}
\rho<1
\quad\Longleftrightarrow\quad
0<\eta<\frac{2\lambda}{\lambda^2+\delta^2}.
\end{equation}
Under this condition, $g_k\to0$ in mean square, while
\begin{equation}
\label{eq:quadratic-exact-init-residual}
E_\infty:=\lim_{k\to\infty}E_k
=\frac{\eta\delta^2\lambda^2x_0^2}
       {2\lambda-\eta(\lambda^2+\delta^2)}.
\end{equation}
Indeed, $e_{t+1}-e_t=-\eta a_{t+1}g_t$ is conditionally centered.
These martingale increments are orthogonal, so for $m>n$,
\[
\mathbb{E}[|e_m-e_n|^2]
=\eta^2\delta^2\sum_{t=n}^{m-1}M_t
\leq\frac{\eta^2\delta^2M_0}{1-\rho}\rho^n
\longrightarrow0.
\]
Thus $e_k$ converges in mean square to some $e_\infty$. Since
$g_k\to0$ in mean square as well,
$x_k=(g_k-e_k)/\lambda$ converges in mean square to
$x_\infty:=-e_\infty/\lambda$. Mean-square convergence gives
$\mathbb E[|e_\infty|^2]=\lim_{k\to\infty}E_k=E_\infty$; hence
\[
\lim_{k\to\infty}M_k=0,
\qquad
\lim_{k\to\infty}G_k
=\mathbb{E}[|\lambda x_\infty|^2]
=E_\infty>0\quad(x_0\neq0).
\]
For $x_0\neq0$, the estimator $g_k$ converges to zero in mean
square, while the true gradient $\nabla f(x_k)$ does not: its
mean-square norm approaches a positive residual. This is \emph{false
convergence}: the estimator signals stationarity even though the
objective has not reached it.

\subsection{Inexact Initialization and Stochastic Transport}
\label{subsec:inexact-init-stochastic-transport}

Let $e_0$ now be arbitrary, while keeping $\delta>0$ and the stable
stepsize \eqref{eq:quadratic-stability}. The recursions
\eqref{eq:quadratic-moment-recursions} still hold, with
$E_0=\mathbb{E}[|e_0|^2]$ and $M_0=\mathbb{E}[|g_0|^2]$. Define
\begin{equation}
\label{eq:quadratic-A}
A:=\frac{\eta^2\delta^2}{1-\rho}
 =\frac{\eta\delta^2}
        {2\lambda-\eta(\lambda^2+\delta^2)}.
\end{equation}
Because $A(1-\rho)=\eta^2\delta^2$, the two moment recursions imply
the conservation law
\begin{equation}
\label{eq:quadratic-conservation}
E_k+A M_k=E_0+A M_0
\qquad\text{for every }k\geq0.
\end{equation}
Since $M_k=\rho^kM_0\to0$, this immediately gives the exact
interaction formula
\begin{equation}
\label{eq:quadratic-interaction}
\boxed{
E_\infty
=E_0+
\frac{\eta\delta^2}
     {2\lambda-\eta(\lambda^2+\delta^2)}M_0
=E_0+A M_0.}
\end{equation}
As above, $G_k\to E_\infty$ as well. The first term is the initial
error that transport never corrects. The second is the total
transport error, governed by the \emph{initial estimator magnitude}
$M_0$, not just by $E_0$.

To see the interaction, write
\[
M_0=\lambda^2x_0^2
    +2\lambda x_0\mathbb{E}[e_0]+E_0.
\]
Bias in $g_0$ can therefore either increase or decrease the
subsequent transport contribution. For the particularly transparent
choice $g_0=\lambda x_0+b_0$, where $b_0$ is centered, independent
of future samples, and $\mathbb{E}[b_0^2]=\sigma^2$, we obtain
\begin{equation}
\label{eq:quadratic-three-contributions}
E_\infty
=\underbrace{\sigma^2}_{\text{initialization}}
 +\underbrace{A\lambda^2x_0^2}_{\text{transport}}
 +\underbrace{A\sigma^2}_{\text{interaction}}.
\end{equation}
The last term is additional transport error caused by the noisy
initial estimator. Initialization quality thus affects both the
persistent offset and the noise accumulated afterward.

\subsection{Summary of the Four Cases}
\label{subsec:four-case-summary}

For the stable quadratic model, the tracking-error limit and the
stationarity limit coincide: $E_\infty=\lim_{k\to\infty}G_k$. The four
cases are summarized below; in the exact-initialization row,
$M_0=\lambda^2x_0^2$.
\begin{center}
\small
\begin{tabular}{@{}lcc@{}}
\toprule
 & Exact transport ($\delta=0$)
 & Stochastic transport ($\delta>0$)\\
\midrule
Exact initialization
 & $E_\infty=0$ (GD)
 & $E_\infty=A\lambda^2x_0^2$\\
Inexact initialization
 & $E_\infty=E_0$ (tilted GD)
 & $E_\infty=E_0+A M_0$\\
\bottomrule
\end{tabular}
\end{center}
The two fundamental error sources are the initial gradient-estimation
error and the errors in transported gradient changes. They interact
because the initial estimate determines the early steps and hence
the subsequent amount of transport noise.
